\documentclass[12pt]{article}
\usepackage{amsmath, amssymb}
\usepackage[margin=1in]{geometry}
\setlength{\parskip}{6pt}
\title{QUBO Formulation: Database Transaction Serializability Verification}
\date{}
\begin{document}
\maketitle

\subsection*{Database Transaction Serializability Verification}

\paragraph{Problem Statement.}
Given a set of $M$ candidate serial transaction orders, each associated with a pre-computed violation cost $c_j \ge 0$ that quantifies the number of dependency constraints violated by order $j$, the goal is to select exactly one candidate order such that the total violation cost is minimized. This corresponds to verifying whether the observed execution is serializable by identifying a candidate with zero violations, or otherwise determining the execution with the minimum number of violations.

\paragraph{Decision Variables.}
We introduce binary variables $x_j \in \{0,1\}$ for each candidate order $j \in \{0, 1, \dots, M-1\}$. The variable $x_j$ indicates the selection of candidate order $j$, where $x_j = 1$ if order $j$ is selected and $x_j = 0$ otherwise.

\paragraph{QUBO Derivation.}
The problem is formulated as a quadratic unconstrained binary optimization (QUBO) by encoding the objective and constraints into a single Hamiltonian $H(\mathbf{x})$.

\textbf{Step 1: Objective Function.}
The objective is to minimize the total violation cost:
\begin{align}
    \min \sum_{j=0}^{M-1} c_j x_j.
\end{align}
This directly contributes to the QUBO Hamiltonian as:
\begin{align}
    H_{\text{obj}}(\mathbf{x}) = \sum_{j=0}^{M-1} c_j x_j.
\end{align}

\textbf{Step 2: Constraint.}
Exactly one candidate order must be selected:
\begin{align}
    \sum_{j=0}^{M-1} x_j = 1.
\end{align}

\textbf{Step 3: Penalty Term.}
We enforce the constraint using a quadratic penalty term with weight $\lambda > 0$:
\begin{align}
    H_{\text{pen}}(\mathbf{x}) = \lambda \left( \sum_{j=0}^{M-1} x_j - 1 \right)^2.
\end{align}
Expanding the square and applying the idempotent property $x_j^2 = x_j$ for binary variables:
\begin{align}
    \left( \sum_{j=0}^{M-1} x_j - 1 \right)^2 &= \sum_{j=0}^{M-1} x_j^2 + \sum_{0 \le j \neq k \le M-1} x_j x_k - 2 \sum_{j=0}^{M-1} x_j + 1 \nonumber \\
    &= \sum_{j=0}^{M-1} x_j + \sum_{0 \le j \neq k \le M-1} x_j x_k - 2 \sum_{j=0}^{M-1} x_j + 1 \nonumber \\
    &= 1 - \sum_{j=0}^{M-1} x_j + \sum_{0 \le j \neq k \le M-1} x_j x_k.
\end{align}
The double sum over distinct indices can be rewritten as twice the sum over ordered pairs $j < k$:
\begin{align}
    \sum_{0 \le j \neq k \le M-1} x_j x_k = 2 \sum_{0 \le j < k \le M-1} x_j x_k.
\end{align}
Substituting this back, the penalty term becomes:
\begin{align}
    H_{\text{pen}}(\mathbf{x}) = \lambda \left( 1 - \sum_{j=0}^{M-1} x_j + 2 \sum_{0 \le j < k \le M-1} x_j x_k \right) \nonumber \\
    = \lambda - \lambda \sum_{j=0}^{M-1} x_j + 2\lambda \sum_{0 \le j < k \le M-1} x_j x_k.
\end{align}

\textbf{Step 4: Combined Hamiltonian.}
The total QUBO Hamiltonian is the sum of the objective and penalty terms:
\begin{align}
    H(\mathbf{x}) &= H_{\text{obj}}(\mathbf{x}) + H_{\text{pen}}(\mathbf{x}) \nonumber \\
    &= \sum_{j=0}^{M-1} c_j x_j + \lambda - \lambda \sum_{j=0}^{M-1} x_j + 2\lambda \sum_{0 \le j < k \le M-1} x_j x_k \nonumber \\
    &= \lambda + \sum_{j=0}^{M-1} (c_j - \lambda) x_j + 2\lambda \sum_{0 \le j < k \le M-1} x_j x_k.
\end{align}

\textbf{Step 5: Q Matrix and Offset.}
Reading off the coefficients from the general form $H(\mathbf{x}) = c + \sum_{j} Q_{j,j} x_j + \sum_{j < k} Q_{j,k} x_j x_k$, we obtain:
\begin{align}
    Q_{j,j} &= c_j - \lambda, \quad \forall j \in \{0, \dots, M-1\}, \nonumber \\
    Q_{j,k} &= 2\lambda, \quad \forall 0 \le j < k \le M-1, \nonumber \\
    c &= \lambda.
\end{align}

\paragraph{Penalty Weight Justification.}
The penalty weight $\lambda$ must be chosen sufficiently large to ensure that any infeasible solution has a strictly higher energy than any feasible solution. The minimum penalty incurred by any infeasible assignment is $\lambda$ (occurring when $\sum x_j = 0$ or $\sum x_j \ge 2$). The maximum possible objective value for a feasible solution is $\max_{j} c_j$. Therefore, choosing
\begin{align}
    \lambda > \max_{j \in \{0, \dots, M-1\}} c_j
\end{align}
guarantees that the energy of any infeasible solution exceeds the energy of any feasible solution, forcing the global minimum to be feasible.

\paragraph{Correctness Argument.}
Minimizing $H(\mathbf{x})$ over $\mathbf{x} \in \{0,1\}^M$ recovers the optimal solution because: (i) for all feasible solutions satisfying $\sum x_j = 1$, the penalty term vanishes and $H(\mathbf{x})$ equals the violation cost $\sum c_j x_j$; (ii) for all infeasible solutions, the penalty term is at least $\lambda$, and by the choice of $\lambda$, the energy of any infeasible solution is strictly greater than the maximum energy of any feasible solution. Consequently, the global minimum of $H(\mathbf{x})$ must be a feasible solution that minimizes the original objective.

\paragraph{Illustrative Example.}
Consider an instance with $M=2$ candidate orders and violation costs $c_0 = 3$, $c_1 = 5$. We choose $\lambda = 10$, which satisfies $\lambda > \max\{3, 5\}$.

The QUBO parameters are:
\begin{align}
    Q_{0,0} &= 3 - 10 = -7, \nonumber \\
    Q_{1,1} &= 5 - 10 = -5, \nonumber \\
    Q_{0,1} &= 2 \times 10 = 20, \nonumber \\
    c &= 10.
\end{align}
The Hamiltonian is:
\begin{align}
    H(x_0, x_1) = 10 - 7x_0 - 5x_1 + 20x_0 x_1.
\end{align}
Evaluating $H$ over all binary assignments:
\begin{align}
    H(0,0) &= 10, \quad \text{(infeasible, penalty } 10\text{)} \nonumber \\
    H(1,0) &= 10 - 7 = 3, \quad \text{(feasible, cost } 3\text{)} \nonumber \\
    H(0,1) &= 10 - 5 = 5, \quad \text{(feasible, cost } 5\text{)} \nonumber \\
    H(1,1) &= 10 - 7 - 5 + 20 = 18. \quad \text{(infeasible, penalty } 10\text{, cost } 8\text{)}
\end{align}
The global minimum is $H(1,0) = 3$, corresponding to selecting candidate order $0$ with minimum violation count $3$.

\end{document}